\documentclass[11pt]{article}
\usepackage[margin=1in]{geometry}
\usepackage[T1]{fontenc}
\usepackage[utf8]{inputenc}
\usepackage{amsmath,amssymb}
\usepackage{booktabs}
\usepackage{longtable}
\usepackage{array}
\usepackage[table]{xcolor}
\usepackage{enumitem}
\usepackage{hyperref}
\usepackage{fancyvrb}
\hypersetup{colorlinks=true, linkcolor=blue, citecolor=blue, urlcolor=blue}

\newcommand{\rtwo}{$R^2$}
\newcommand{\Rsq}{R^{2}}
\newcommand{\HSL}{\textsc{HyperSymLoop}}
\newcommand{\EHD}{\textsc{EHSDeFi}}
\newcommand{\open}{\textcolor{red!60!black}{\textbf{[OPEN --- needs code/data investigation]}}}
\newcommand{\stale}{\textcolor{orange!80!black}{\textbf{[STALE --- reconciliation only]}}}
\newcommand{\textfix}{\textcolor{green!40!black}{\textbf{[TEXT --- pure writing fix]}}}
\newcommand{\done}{\textcolor{blue!60!black}{\textbf{[ALREADY RESOLVED --- doc propagation only]}}}

\newcommand{\loc}[1]{\textit{Location:} #1\par}
\newenvironment{evidence}
  {\begin{quote}\small\ttfamily}
  {\end{quote}}

\title{\bf Wrong-Section Report\\\large Thirteen Confirmed Discrepancies, Cross-Referenced to Source}
\author{}
\date{Compiled 2026-07-31}

\begin{document}
\maketitle
\thispagestyle{empty}

\begin{abstract}
This report locates, quotes, and sub-labels every wrong or internally
inconsistent section identified in the three manuscript files
(\texttt{jmlr\_paper\_main.tex}, \texttt{supp\_benchmark\_report.tex},
\texttt{supp\_routing\_improvements.tex}). Each numbered section below
corresponds one-to-one with the issue of the same number in the audit
fix list, gives the exact file/label/line location of the conflicting
statements, quotes the conflicting text verbatim, and states the category
of fix required.
\end{abstract}

\tableofcontents
\clearpage

% =======================================================================
\section{Nguyen-12 Table Success-Count Arithmetic}
\label{issue:1}
% =======================================================================
\stale\ (partial arithmetic fix) \open\ (per-equation values unverified)

\loc{\texttt{jmlr\_paper\_main.tex}, \S\texttt{subsec:nguyen12-benchmark-standard-sr-suite}
``Nguyen-12 Benchmark (Standard SR Suite)'', \texttt{tab:nguyen12}, l.\,2432--2520.}

\subsection*{What is wrong}
The table's printed \emph{Success} row and its own per-equation cells
disagree with each other under the table's own stated criterion
(``Bold: extrap \rtwo\ $\geq 0.9999$''):

\begin{evidence}
Success (\rtwo$\ge$0.9999) \quad P: 10/12 (83.3\%) \quad H: 11/12 (91.7\%) \quad N: 0/12 (0\%)
\end{evidence}

A hand recount of the twelve bolded per-equation extrapolation cells gives:
\textbf{PySR-only (P): 8/12}, not 10/12 (N-3 and N-7 are not $\geq0.9999$ but
are not counted as passes, and no other cell was miscounted --- the printed
total simply does not equal the number of bolded cells). \textbf{HypatiaX
(H): 7/12}, not 11/12 --- two cells (N-3 $=0.9976$, N-10 $=0.9997$) are
typeset in \textbf{bold} despite being below the table's own $\geq0.9999$
threshold, and removing them from the bolded/passing count leaves 7, not 11.
A third, independently conflicting figure appears in the caption itself:

\begin{evidence}
Strict recovery ($\ge 0.9999$) yields 4/12 (33.3\%); the 91.7\,\% figure uses
the 4-decimal rounding convention.
\end{evidence}

Three different denominators for the same nominal quantity --- 10/12 (row),
11/12 (row, also restated in prose at l.\,220, 2503), and 4/12 (caption) ---
appear in the same table environment and do not reconcile against each other
or against the twelve visible per-equation cells.

\subsection*{Fix required}
The Success row is fixable by hand re-tally against the table's own printed
cells. However, per the fix-list's caution, the per-equation \rtwo\ values
themselves have not been independently re-verified against the raw result
JSON, so a complete fix should regenerate both the per-equation cells and
the summary row from source, not merely re-tally the printed numbers.

% =======================================================================
\section{Supp.\ Table 4 (30/30) vs.\ Abstract/Conclusion (27/30) on \HSL{} (M4)}
\label{issue:2}
% =======================================================================
\stale\ --- classic same-run desync, no code investigation needed

\loc{\texttt{supp\_benchmark\_report.tex}, \texttt{tab:overall}
(``Six-Method Aggregate Performance,'' l.\,394--423) vs.\ its own Abstract
(l.\,140--141).}

\subsection*{What is wrong}
The Abstract states:
\begin{evidence}
\HSL{} achieves 100\% recovery at all \emph{noisy} conditions ($\sigma > 0$)
and 90.0\% (27/30 equations) under the strict noiseless protocol
($\Rsq > 0.9999$).
\end{evidence}
But \texttt{tab:overall}, in the same document, prints:
\begin{evidence}
\HSL\,(M4)\quad Core\quad 30/30 (v2)$^\dagger$\quad 0.999+\quad 11.4\,s
\end{evidence}
with a footnote reading ``\HSL{} v1 showed 26/30 due to the Newton
measurement bug; v2 achieves 30/30.'' The table's 30/30 and the abstract's
27/30 are both presented as the current (v2, corrected) figure for the same
method on the same noiseless protocol, in the same document.

\subsection*{Fix required}
Regenerate \texttt{tab:overall}'s \HSL{} row from whichever run actually
produced the 27/30 figure quoted in the abstract, and confirm both numbers
come from the same run before republishing either.

% =======================================================================
\section{\EHD{} (M3) Runtime: 20.2\,s vs.\ 841.4\,s}
\label{issue:3}
% =======================================================================
\open\ --- needs live code/data investigation

\loc{\texttt{supp\_benchmark\_report.tex}, \texttt{tab:overall} l.\,413 vs.\
\texttt{tab:time\_noise} (``Sweep Results: Noise Robustness'') l.\,571--584.}

\subsection*{What is wrong}
\texttt{tab:overall} reports \EHD{}'s (M3) average runtime as:
\begin{evidence}
\EHD\,(M3)\quad Core\quad 29/30 (96.7\%)\quad 0.999993\quad 20.2\,s
\end{evidence}
\texttt{tab:time\_noise}, covering the same method at $\sigma=0\%$ (the same
noiseless condition), reports:
\begin{evidence}
0\% \quad 841.4 \quad 11.1 \quad \textbf{75.8$\times$}
\end{evidence}
The abstract's own claim that ``\EHD{} offers exact symbolic
interpretability; \HSL{} is $1.5$--$76\times$ faster depending on noise
level'' corroborates the $841.4$\,s figure (75.8$\times$ vs.\ M4's 11.1\,s),
not the 20.2\,s figure.

\subsection*{Fix required}
Trace which script/run produced 20.2\,s. If it reflects different
hardware/config, that must be stated explicitly; if it is simply wrong, it
must be regenerated from the same raw timing logs that produced
\texttt{tab:time\_noise}.

% =======================================================================
\section{CI / SD Statistical Incompatibility}
\label{issue:4}
% =======================================================================
\open\ --- needs live code/data investigation

\loc{\texttt{jmlr\_paper\_main.tex}, \S\texttt{subsec:fivesystem-comparative-analysis-feynman-}
``Five-System Comparative Analysis (Feynman Core-15),'' \texttt{tab:five\_systems\_full}
and \texttt{thm:five\_system\_hierarchy}, l.\,1213--1259 vs.\
\texttt{supp\_benchmark\_report.tex}, Appendix G ``Statistical Test
Details,'' \texttt{app:statistical\_tests}, l.\,1423--1471.}

\subsection*{What is wrong}
The main paper's Result box states:
\begin{evidence}
neural networks exhibit mean error 1231\% (95\% CI: [1087\%, 1456\%], $n=13$)
\end{evidence}
That gives a CI half-width of $\approx 185$ percentage points around a mean
of 1231\%. Appendix G, describing the same neural-network distribution,
states:
\begin{evidence}
Neural Network: Mean = 1231\%, Median = 86.7\%, SD = 1842\%,
Range = [12.3, 5847\%] (CORRECTED)
\end{evidence}
A standard 95\% CI on the mean with $n=13$ and $\mathrm{SD}=1842\%$ has a
half-width of roughly $t_{0.975,12}\cdot 1842/\sqrt{13} \approx 2.18 \times
511 \approx 1114$ percentage points --- almost six times narrower a CI is
reported in the main text than the disclosed SD supports.

\subsection*{Fix required}
Re-run whatever CI computation produced Table 1's figure and check it
against Appendix G's SD; the CI was very likely computed with the wrong SD,
wrong $n$, or an undisclosed non-standard method.

% =======================================================================
\section{``68 of 74'' Cited Without Seed/Split}
\label{issue:5}
% =======================================================================
\textfix\ --- pure writing/citation fix, underlying data is consistent

\loc{\texttt{jmlr\_paper\_main.tex}, Abstract l.\,222--223 and
\S\texttt{subsec:fivestage-routing-architecture-overview} l.\,1011,
referencing Table~\texttt{tab:timing\_full} (\S\texttt{subsec:runtime-analysis}).}

\subsection*{What is wrong}
Both the abstract and \S\texttt{subsec:fivestage-routing-architecture-overview}
cite ``68 of 74'' LLM-routed tasks with no seed or split attached:
\begin{evidence}
A previously reported $1.73\times$ median speedup of HypatiaX over
neural-network inference on LLM-routed cases (68 of 74) does not reproduce\ldots
\end{evidence}
\begin{evidence}
Stage~3 (optional): LLM proposes candidate expressions; on the v3.0 DeFi
benchmark, 68 of 74 tasks are routed to this path.
\end{evidence}
But Table~\texttt{tab:timing\_full} shows \emph{two different} rows with exactly
68/74 LLM-routed tasks --- \texttt{v3c seed123} (68/74) and \texttt{PCA
seed42} (68/74) --- which are different seeds under different splits with
different absolute runtimes. ``68 of 74,'' as written, does not identify
which of the two.

\subsection*{Fix required}
Attach the specific seed and split (e.g.\ ``68 of 74 under v3c seed 123'' or
``\ldots PCA seed 42'') wherever ``68 of 74'' is cited. No table value needs
to change.

% =======================================================================
\section{``v3c'' vs.\ ``PCA-directed'' Never Defined as Two Variants}
\label{issue:6}
% =======================================================================
\textfix\ --- pure writing fix, both variants' numbers are real

\loc{\texttt{jmlr\_paper\_main.tex}, \S\texttt{subsec:extrapolation-protocol}
(\texttt{sec:split}, l.\,678--709) vs.\ Table~\texttt{tab:timing\_full} /
Table~\texttt{tab:timing\_llm\_routed\_full} captions (l.\,1630--1636).}

\subsection*{What is wrong}
\S\texttt{subsec:extrapolation-protocol} (``Extrapolation Protocol'') defines
exactly \emph{one} split methodology:
\begin{evidence}
To rigorously assess out-of-distribution generalisation, we employ a
\textbf{PCA-directed aggressive extrapolation split}: \ldots
\end{evidence}
No ``v3c'' variant is introduced or distinguished anywhere in this section.
Yet Table~\texttt{tab:timing\_full}'s caption (and every runtime table
downstream) treats \texttt{v3c} and \texttt{PCA} as two distinct,
sibling splits:
\begin{evidence}
Rows are grouped by split (v3c = aggressive 40/60; PCA = PCA-directed 40/60)\ldots
\end{evidence}
A reader following \S\texttt{subsec:extrapolation-protocol} alone has no way to
learn that ``v3c'' and ``PCA-directed'' name two different, non-identical
protocols rather than being two names for the same one.

\subsection*{Fix required}
Add an explicit definition of both variants (what specifically distinguishes
\texttt{v3c}'s ``aggressive'' routing from the PCA-directed routing already
described) to \S\texttt{subsec:extrapolation-protocol}, since the underlying
runs for both are real.

% =======================================================================
\section{Remark on Difficulty-Tier Counts: $24/27/20 \to 24/29/21$ vs.\ ``+3 Hard''}
\label{issue:7}
% =======================================================================
\textfix\ --- pure text/arithmetic fix

\loc{\texttt{jmlr\_paper\_main.tex}, \S\texttt{subsec:difficulty-classification}
unlabeled remark following the Easy/Medium/Hard definitions, l.\,656--660.}

\subsection*{What is wrong}
The tier definitions immediately above state Easy $n=24$, Medium $n=29$,
Hard $n=21$ (totalling 74). The remark directly below reads:
\begin{evidence}
An earlier version of this benchmark comprised 71 tasks (Easy=24, Medium=27,
Hard=20). The current benchmark (v3.0, 74 tasks) adds three additional hard
cases in multi-asset risk: correlated portfolio VaR, component ES, and
multi-collateral LTV.
\end{evidence}
$71 + 3 = 74$ is correct in total, but $24/27/20 \to 24/29/21$ is a change
of $+0/+2/+1$, not $+0/+0/+3$. The remark's claim that all three additions
are Hard-tier is arithmetically inconsistent with the tier counts stated one
paragraph above it.

\subsection*{Fix required}
Either the remark's ``adds three additional hard cases'' is wrong (should
read something like ``+2 Medium, +1 Hard''), or the Medium/Hard counts in
\S\texttt{subsec:difficulty-classification} are wrong. Purely a text
correction; no code or data investigation implicated.

% =======================================================================
\section{Figure 3 Caption ``DeFi AMM (Hard) $=-78.04$'' Untraceable, Plus Impermanent-Loss Tier Clash}
\label{issue:8}
% =======================================================================
\open\ --- confirmed known-bad figure/caption, fix source already identified

\loc{\texttt{jmlr\_paper\_main.tex}, Figure~\texttt{fig:r2\_heatmap\_clipped}
caption (\texttt{fig18\_r2\_heatmap\_improved}, l.\,1287--1300) vs.\
Table~\texttt{tab:llm\_ablation} (l.\,1871--1897) vs.\ the Medium-tier
definition (l.\,632--634).}

\subsection*{What is wrong}
The Figure~3 caption reads:
\begin{evidence}
HypatiaX is not uniformly better than PySR-only: it improves on Chemistry
(Easy: $0.90$ vs.\ $-12.56$) but produces $-\infty$ --- a genuine evaluation
crash, not merely a poor fit --- on DeFi AMM (Hard) and Physics (Hard),
where PySR-only instead returns a finite (if poor, $-78.04$, or perfect,
$1.00$) value.
\end{evidence}
Cross-checking against Table~\texttt{tab:llm\_ablation} (explicitly stated to
underlie this figure): no ``DeFi AMM'' equation shows PySR-only $=-78.04$
at any range (Constant Product P: $0.9982$--$0.9996$; Impermanent Loss P:
$0.7703$ near, $-2.7832$ med, $-157.0776$ far; Price Impact P: all
$\approx1.0$). The $-\infty$ HypatiaX crash described for the DeFi AMM row
also does not appear in the DeFi AMM rows of Table~\texttt{tab:llm\_ablation}
(all DeFi AMM H values are finite, if very large-magnitude negative); the
only $-\infty$ H value in the whole table belongs to \textbf{Gravitational
Force}, which is Physics, not DeFi AMM. Separately, the caption implies a
DeFi AMM task sits in the \textbf{Hard} tier, but the difficulty-tier
definition explicitly lists Impermanent Loss --- the DeFi AMM equation
actually exhibiting catastrophic values in Table~\texttt{tab:llm\_ablation}
--- as a \textbf{Medium}-tier example:
\begin{evidence}
Medium ($n=29$): nonlinear algebraic or exponential relationships, such as
impermanent loss, Uniswap V3 liquidity range, constant-product price
impact\ldots
\end{evidence}

\subsection*{Fix required}
Regenerate Figure~3 from the same \texttt{\_merged.json} data used to
rebuild Table~\texttt{tab:llm\\_ablation}, rather than whatever process produced
the caption's $-78.04$/DeFi-AMM-Hard claims, and resolve whether an
Impermanent-Loss-class task belongs to Medium or Hard consistently across
the tier definition and the figure.

% =======================================================================
\section{Uncorrected Mean \rtwo\ ($+0.8721$) Despite Disclosed $-141{,}000$ Masked Failure}
\label{issue:9}
% =======================================================================
\open\ --- pending, not yet executed against a full live re-run

\loc{\texttt{jmlr\_paper\_main.tex}, \S\texttt{subsec:overall-extrapolation-performance-defi-v}
\texttt{tab:main\_results}, l.\,1261--1285, and the Abstract's masking-disclosure
footnote, l.\,209--215.}

\subsection*{What is wrong}
The abstract discloses that HypatiaX's own success accounting silently masks
catastrophic sub-method failures:
\begin{evidence}
The underlying computation is not free of catastrophic failures: 22 of the
74 tasks route to a sub-method (typically the LLM proposal) that itself
failed catastrophically (e.g.\ \texttt{pure\_llm.test\_r2} $\approx
-141{,}000$ on Loan-to-Value), but the hybrid system's own accounting
reports \texttt{success=True}, $\Rsq \approx 1.0$ regardless. The
catastrophic failure is masked, not eliminated.
\end{evidence}
Yet \texttt{tab:main\_results}, whose own caption acknowledges the
$90.5\%$ pass-rate figure is uncorrected for this exact bug, still reports
HypatiaX's \textbf{Mean \rtwo} as $+0.8721$ --- computed, like the pass
rate, from the same masked \texttt{success}/\rtwo\ fields that report
$\approx1.0$ for the 22 tasks whose true sub-method result is catastrophic:
\begin{evidence}
HypatiaX \quad 1.0000 \quad $+0.8721$ \quad \textbf{90.5} \quad \textbf{90.5} \quad \textbf{0}
\end{evidence}
The caption corrects the pass-rate columns (``the corrected overall
near-perfect rate is $60.8\%$'') but is silent on the Mean \rtwo\ column,
which is left standing as if unaffected by the same bug.

\subsection*{Fix required}
Recompute Mean \rtwo\ from the corrected (decision-attribution-fixed) full
live run described in \S\texttt{sec:hybrid-attribution-bug}; this is
explicitly stated elsewhere as not yet executed, so the corrected number
cannot be hand-estimated from the single disclosed $-141{,}000$ outlier
alone.

% =======================================================================
\section{Timeout Contradiction: 1100\,s vs.\ 900\,s; ``300\,s Limit'' but a 27{,}574\,s Hang}
\label{issue:10}
% =======================================================================
\open\ (enforcement bug, second half) / \stale\ (config bleed, first half)

\loc{\texttt{supp\_benchmark\_report.tex}, Sweep hyperparameters
(\texttt{tab:sweeps}, l.\,294--311) vs.\ Reproducibility Statement, Phase~3
reproduction command, l.\,988--1004.}

\subsection*{What is wrong}
\texttt{tab:sweeps}' hyperparameter note states:
\begin{evidence}
Hyperparameters: NN seeds = 5, random seed 42, method timeout = 900\,s,
PySR timeout = 1100\,s, threshold (noisy) = 0.995, threshold (noiseless) =
0.999999. All 30 equations completed with zero timeouts across all 11
conditions.
\end{evidence}
The Phase~3 reproduction command block, describing (per its own caption)
the run underlying \texttt{tab:overall}, instead shows:
\begin{evidence}
python run\_comparative\_suite\_benchmark\_v2.py \textbackslash\\
\ \ --noiseless --threshold 0.9999\ \ \ \ \ \ \ \textbackslash\\
\ \ --nn-seeds 3 --samples 200\ \ \ \ \ \ \ \ \ \ \ \textbackslash\\
\ \ --method-timeout 900\ \ \ \ \ \ \ \ \ \ \ \ \ \ \ \ \ \ \ \ \ \textbackslash\\
\ \ --pysr-timeout 900
\end{evidence}
i.e.\ 900\,s for \emph{both} timeouts, not 1100\,s for PySR as
\texttt{tab:sweeps} states. Immediately below, the same reproducibility
note directly contradicts the ``zero timeouts'' claim:
\begin{evidence}
Note on timeout upgrade: \texttt{--method-timeout} was set to 300\,s for
tests 1--18 and upgraded to 900\,s for tests 19--30. Arrhenius (test~4) ran
under the 300\,s limit and hung for 27574\,s before the Julia process was
killed\ldots Nernst (test~6) also timed out under the 300\,s default;
counted as failure.
\end{evidence}
A stated ``300\,s limit'' allowing a process to run for 27{,}574\,s (over
90$\times$ the configured limit) before manual kill indicates the timeout
was not actually enforced for that run.

\subsection*{Fix required}
First half (1100\,s vs.\ 900\,s): reads as two different phases' configs
bleeding into one write-up --- reconcile which value is correct for
\texttt{tab:sweeps} and state it consistently. Second half (300\,s limit,
27{,}574\,s hang): the timeout \emph{enforcement} itself needs checking in
the harness code (or was silently disabled for that run) --- this is a
code-level question, not just a documentation fix.

% =======================================================================
\section{Model/Temperature Mismatch Across Documents}
\label{issue:11}
% =======================================================================
\done\ --- root cause already found; only stale doc entries remain

\loc{\texttt{supp\_benchmark\_report.tex}, reproducibility tables l.\,1045--1046
and l.\,1243/1261, and \texttt{supp\_routing\_improvements.tex}, ``Hardware
and Software'' l.\,756, vs.\ \texttt{jmlr\_paper\_main.tex}, Nguyen-12
reproducibility note, l.\,2606--2637.}

\subsection*{What is wrong}
\texttt{jmlr\_paper\_main.tex} already contains the corrected finding: the
model that actually produced the Nguyen-12 results is
\texttt{claude-sonnet-5} at \texttt{temperature=0.25}, via \texttt{hypatia.py}
--- confirmed as a fourth, distinct hardcoded string, with the
\texttt{LLM\_MODEL} environment variable and \texttt{repro.yaml}'s
\texttt{llm\_model} field both confirmed to be dead code that nothing reads.
Three stale locations in the supplementary files still report the
superseded value:
\begin{evidence}
supp\_benchmark\_report.tex:1045 \quad Anthropic SDK \quad 0.73.0
(\texttt{claude-sonnet-4-20250514}, temperature~0.0)
\end{evidence}
\begin{evidence}
supp\_benchmark\_report.tex:1243, 1261 \quad Model: claude-sonnet-4-20250514
\quad LLM temperature: 0.3
\end{evidence}
\begin{evidence}
supp\_routing\_improvements.tex:756 \quad Anthropic SDK \quad 0.73.0 \quad
\texttt{claude-sonnet-4-20250514}, temp~0.0
\end{evidence}

\subsection*{Fix required}
Overwrite all three stale entries with the confirmed value
(\texttt{claude-sonnet-5}, \texttt{temperature=0.25}). No further
investigation needed --- this is doc propagation only.

% =======================================================================
\section{Rounding Slippage: $1.21/0.36 = 3.36\times$, Printed $3.40\times$}
\label{issue:12}
% =======================================================================
\open\ --- narrow, mechanical script fix + regen

\loc{\texttt{jmlr\_paper\_main.tex}, \S\texttt{subsec:runtime-analysis}
Table~\texttt{tab:timing\_full} pooled PCA row, l.\,1636.}

\subsection*{What is wrong}
The pooled PCA row of Table~\texttt{tab:timing\_full} reads:
\begin{evidence}
\textbf{PCA ALL SEEDS (pooled)} \quad \textbf{370} \quad \textbf{2.41 / 0.18}
\quad \textbf{0.36 / 0.35} \quad \textbf{1.21 / 0.89} \quad \textbf{94/370}
\quad \textbf{3.40$\times$ slower}
\end{evidence}
Dividing the printed Hybrid mean (1.21\,s) by the printed Neural MLP mean
(0.36\,s) gives $1.21/0.36 \approx 3.361\times$, not the printed
$3.40\times$. (By contrast every v3c-split row's printed speedup is
internally consistent with its own printed mean columns.)

\subsection*{Fix required}
\texttt{generate\_table1.py} (already named in the paper as the
regeneration script) is almost certainly dividing pre-rounded display
values instead of full-precision ones; needs a one-line fix plus a full
table regeneration.

% =======================================================================
\section{Near-Duplicate Task Names Across Difficulty Tiers}
\label{issue:13}
% =======================================================================
\textfix\ --- naming-clarity issue, not a data bug

\loc{\texttt{jmlr\_paper\_main.tex}, \S\texttt{sec:hybrid-attribution-bug}
fabricated-success task list, l.\,1540--1550, vs.\
\S\texttt{subsec:lending-protocols} ``Lending Protocols'' task inventory,
l.\,2960--2977.}

\subsection*{What is wrong}
The Lending Protocols task inventory lists two visually near-identical
pairs of tasks as distinct benchmark items:
\begin{evidence}
Liquidation Price Long and Liquidation Price Short: prices at which a
position is liquidated.
\ldots
Liquidation price for leveraged long and leveraged short: extended formulas
accounting for funding rates.
\end{evidence}
Both pairs then appear, split across different difficulty tiers, in the
fabricated-success task list:
\begin{evidence}
Medium (10): \ldots Liquidation Price Long, Liquidation Price Short, \ldots
\end{evidence}
\begin{evidence}
Hard (9): \ldots Liquidation price for leveraged long, Liquidation price
for leveraged short, \ldots
\end{evidence}
A reader skimming either list alone could easily mistake ``Liquidation
Price Long'' (Medium) for ``Liquidation price for leveraged long'' (Hard),
or assume a duplicate/copy-paste error, when in fact these name four
formulas that are genuinely distinct (the leveraged variants explicitly
account for funding rates).

\subsection*{Fix required}
Rename one pair for clarity (e.g.\ ``Liquidation Price, Basic Long/Short''
vs.\ ``Liquidation Price, Leveraged Long/Short (funding-adjusted)'') so the
two are visually distinguishable wherever they are listed. No numeric or
code change needed.

\clearpage
% =======================================================================
\section*{Summary Table}
% =======================================================================
\begin{longtable}{p{1.2cm}p{6.7cm}p{2.2cm}p{2.8cm}}
\toprule
\textbf{\#} & \textbf{Issue} & \textbf{Fix category} & \textbf{Primary files} \\
\midrule
\endhead
\ref{issue:1} & Nguyen-12 success-count arithmetic (10/12, 7-actual, 4/12) & Stale + Open & jmlr\_paper\_main \\
\ref{issue:2} & Supp Table 4 (30/30) vs.\ Abstract (27/30), M4 & Stale & supp\_benchmark\_report \\
\ref{issue:3} & \EHD{} runtime 20.2\,s vs.\ 841.4\,s & Open & supp\_benchmark\_report \\
\ref{issue:4} & CI/SD statistical incompatibility & Open & jmlr\_paper\_main / supp\_benchmark\_report \\
\ref{issue:5} & ``68 of 74'' without seed/split & Text & jmlr\_paper\_main \\
\ref{issue:6} & v3c vs.\ PCA-directed undefined & Text & jmlr\_paper\_main \\
\ref{issue:7} & Remark 9 tier arithmetic ($+3$ Hard) & Text & jmlr\_paper\_main \\
\ref{issue:8} & Fig.\ 3 caption $-78.04$ untraceable + tier clash & Open (source ID'd) & jmlr\_paper\_main \\
\ref{issue:9} & Uncorrected Mean \rtwo\ despite masked failure & Open (pending re-run) & jmlr\_paper\_main \\
\ref{issue:10} & Timeout contradiction (1100 vs.\ 900\,s; 300\,s limit / 27{,}574\,s hang) & Open + Stale & supp\_benchmark\_report \\
\ref{issue:11} & Model/temperature mismatch & Already resolved (doc only) & supp\_benchmark\_report / supp\_routing\_improvements \\
\ref{issue:12} & Rounding slippage 3.36$\times$ vs.\ 3.40$\times$ & Open (mechanical) & jmlr\_paper\_main \\
\ref{issue:13} & Near-duplicate task names & Text & jmlr\_paper\_main \\
\bottomrule
\end{longtable}

\end{document}
